\section{The complete tame elliptic criterion}\label{sec:tame-full-torsion}
This section proves Theorems~\ref{thm:tame-level-one}--\ref{thm:tame-full-level} and Corollary~\ref{cor:tame-forcing-level}, including actual elliptic realizations at every allowed absolute ramification index.

\subsection{Inertia and the good formal group}
We first identify the good special fibre and the restriction imposed by a tame tangent character.

\begin{lemma}\label{tf-lem:inertia}
Let $p$ be odd, and let $E/K$ be potentially good with nontrivial tame semistability defect $m$ and full rational $p$-torsion. Its good special fibre is supersingular. If $p\ge5$, every nontrivial potentially good defect is tame and belongs to $\{2,3,4,6\}$; for $p=3$ a nontrivial tame defect belongs to $\{2,4\}$.

After completed maximal unramified base change, let $L/K$ be the minimal good extension. There is an integral formal parameter $t$ and a residue class $r\in(\Z/m\Z)^\times$ such that every nonzero rational descended $p$-power torsion point satisfies
\begin{equation}\label{tf-eq:congruence}
 v_L(t(P))\equiv r\pmod m.
\end{equation}
\end{lemma}

\begin{proof}
Completed maximal unramified base change preserves the absolute index and the inertia image. We may therefore assume the residue field $k$ is algebraically closed. We use the good-reduction criterion of \citet[Theorem~1 and Corollary~4, pp.~493--494]{SerreTate1968}. The finite inertia image acts faithfully on the good special elliptic curve $E_0$: specialization identifies prime-to-$p$ Tate modules, and an automorphism acting trivially on every $\ell$-power torsion point is the identity. Indeed a nonzero elliptic endomorphism is finite and cannot kill an infinite Tate module. The good-reduction criterion identifies the kernel with the good-reduction inertia kernel.

For $p\ge5$, a short Weierstrass equation $y^2=x^3+Ax+B$ identifies every origin-preserving automorphism with
$(x,y)\mapsto(u^2x,u^3y)$, where $u^4A=A$ and $u^6B=B$. Its group is cyclic of order $2$, $4$ or $6$, and all subgroups are tame. In characteristic three, complete the square and write
$y^2=x^3+a_2x^2+a_4x+a_6$. An automorphism has
$x\mapsto u^2x+r$, $y\mapsto u^3y$. If $a_2\ne0$, comparison of the quadratic coefficients gives $u^2=1$. If $a_2=0$, nonsingularity gives $a_4\ne0$, and comparison of linear coefficients gives $u^4=1$. A finite automorphism of order prime to $p$ with tangent multiplier one is the identity: its first nonidentity term $cT^j$ would become $hcT^j$ on taking an $h$th iterate, which cannot vanish when $p\nmid h$. Thus a tame subgroup injects into the tangent multipliers. This proves the asserted lists and the faithfulness of the tame tangent character.

Suppose $E_0$ were ordinary. The finite flat good-model group scheme $E[p]$ shows that reduction of geometric generic $p$-torsion surjects onto $E_0[p](k)$: each special point lies in the closure of generic points. All generic points are rational over $L$, so the descent automorphism acts trivially on these special points. Its action on the height-one \emph{\'etale} $p$-divisible group is a finite-order element of $\Z_p^*$ trivial modulo $p$. The group $1+p\Z_p$ is torsion-free for odd $p$, so this action is trivial. The endomorphism $\gamma-1$ kills the entire \emph{\'etale} Tate module and hence is zero. This contradicts nontrivial inertia. Therefore $E_0$ is supersingular and all geometric $p$-power torsion lies in the good formal group.

The tame extension $L/K$ is cyclic of degree $m$. Choose a uniformizer $\pi$ and generator $\sigma$ with $\sigma\pi=\xi\pi$, where $\xi\in K$ is a primitive $m$th root. This follows from the tame Kummer description over a strictly henselian DVR; units have $m$th roots by Hensel's lemma. Write the faithful tangent multiplier as $\lambda=\xi^r$, with $\gcd(r,m)=1$, using the action convention on functions. For any formal parameter $t_0$, the average
\[
 t=\frac1m\sum_{i=0}^{m-1}\lambda^{-i}\sigma^i(t_0)
\]
is again a parameter and satisfies $\sigma(t)=\lambda t$. Its special linear coefficient is a unit. Evaluation on a descended torsion section gives
$\sigma(t(P))=\lambda t(P)$. If $q=v_L(t(P))$, the residue of $t(P)/\pi^q$ is fixed by inertia, so $\xi^q=\lambda$ in $k$. This gives~\eqref{tf-eq:congruence}.
\end{proof}

\subsection{The smallest and largest division valuations}
The next calculation retains both extrema of exact-order torsion and includes every equality case.

\begin{lemma}\label{tf-lem:extrema}
Let $F$ be a one-dimensional formal group of special height two over a complete mixed-characteristic DVR, with full rational $p^n$-torsion. The residue field may be imperfect. Normalize the valuation to $\Z$, put $E=v(p)$, and let $a\ge b>0$ be the largest and smallest nonzero $p$-torsion valuations. Then
\begin{equation}\label{tf-eq:mass}
 M:=\frac{E}{p-1}=a+pb.
\end{equation}
If $b_n,A_n$ are the minimum and maximum exact $p^n$-torsion valuations, then
\begin{equation}\label{tf-eq:extrema}
 b_n=\frac b{p^{2n-2}},\qquad
 A_n=\frac{M-p^{c+1}b}{p^{2n-2-c}},\qquad
 c=\min\left\{n-1,\left\lfloor\log_p\frac{M}{(p+1)b}\right\rfloor\right\}.
\end{equation}
\end{lemma}

\begin{proof}
The positive-valuation filtration on $F[p]$ is a filtration of a two-dimensional $\F_p$-space. Indeed $F(X,Y)=X+Y$ plus terms divisible by $XY$, so addition has valuation at least the minimum, with equality at unequal valuations; multiplication by a $p$-adic unit preserves valuation. A strictly nested filtration has at most one nonzero proper subspace. Consequently, if $a>b$, exactly $p-1$ nonzero points have value $a$ and $p^2-p$ have value $b$. If $a=b$, all $p^2-1$ points have that value.

Weierstrass preparation of $[p](T)$ has distinguished degree $p^2$. Comparing its derivative at zero with the product of its nonzero roots gives~\eqref{tf-eq:mass}. The Newton polygon has vertices
\[
 (1,E),\qquad(p,H),\qquad(p^2,0),\qquad H=p(p-1)b,
\]
with the middle point collinear in the one-slope case. This follows from the root valuations and multiplicities and requires no special coordinate normal form.

Let $P\ne0$ be a torsion point of valuation $q\le a$. The lower polygon for $[p](T)-t(P)$ governing positive root valuations is the hull of
\[
 (0,q),\qquad(p,H),\qquad(p^2,0).
\]
The point $(1,E)$ lies above this hull because $E>q$; all old coefficient points lie above the old polygon, which lies above the proposed new hull. Coefficients beyond degree $p^2$ have nonnegative valuations. Thus the division values are
\begin{equation}\label{tf-eq:division}
\begin{cases}
 q/p^2\quad\text{for all }p^2\text{ preimages},&q\le p^2b,\\
 (q-H)/p\quad\text{for }p\text{ preimages, and }b
 \text{ for }p^2-p\text{ preimages},&q>p^2b.
\end{cases}
\end{equation}
At equality both descriptions coincide. This formula is used only for nonzero torsion with $q\le a$. Multiplication by $p$ strictly increases a positive nonzero formal valuation, so every higher-order torsion valuation is smaller than $a$ and the required range is preserved.

The minimum branch starting at $b$ divides by $p^2$ at every step, giving the first formula in~\eqref{tf-eq:extrema}. The largest division value is the continuous increasing function
\[
 f(q)=\begin{cases}q/p^2,&q\le p^2b,\\
 (q-p(p-1)b)/p,&q\ge p^2b.
 \end{cases}
\]
Hence $A_n=f^{n-1}(a)$. On the upper branch, after $j-1$ divisions the value is $M/p^{j-1}-pb$. The integer $c$ in~\eqref{tf-eq:extrema} counts the upper steps until the value enters $[b,p^2b]$, or until the chain ends. The remaining steps divide by $p^2$, and give the second displayed formula. If $M/((p+1)b)$ is an exact power of $p$, the two adjacent branch choices agree at their common endpoint. If $a=b$, then $c=0$ and both extrema are $b/p^{2n-2}$.
\end{proof}

\begin{proposition}\label{tf-prop:necessity}
Every curve in Theorem~\ref{thm:tame-full-level} satisfies~\eqref{eq:tame-full-criterion}. Full rational $p$-torsion with nontrivial tame defect also forces $m\mid p+1$ and $p-1\mid e$.
\end{proposition}

\begin{proof}
The Weil pairing gives $\mu_{p^n}\subset K$, and
\[
 v_K(\zeta_{p^n}-1)=\frac{e}{p^{n-1}(p-1)}\in\Z.
\]
In particular $p-1\mid e$. On the minimal good field, $E=me$ and
$M=me/(p-1)$. Lemma~\ref{tf-lem:inertia} gives $a\equiv b\equiv r\pmod m$ with $r$ invertible. Since $m\mid M=a+pb$, we obtain $(p+1)r\equiv0\pmod m$, hence $m\mid p+1$.

Put $s=v_p(e/(p-1))=v_p(M)$. Integrality of $b_n$ in Lemma~\ref{tf-lem:extrema} gives
\[
 p^{2n-2}\mid b,\qquad b\ge p^{2n-2}.
\]
Integrality of $A_n$ gives $p^{2n-2-c}\mid M-p^{c+1}b$. The second term is divisible by $p^{2n+c-1}$, so
\[
 s\ge2n-2-c\ge n-1.
\]
The definition of $c$, including its capped case, gives $M\ge(p+1)bp^c$. Therefore
\[
 M\ge(p+1)p^{2n-2+\max(0,2n-2-s)}
       =(p+1)p^{\max(2n-2,4n-4-s)}.
\]
Multiplication by $(p-1)/m$ proves~\eqref{eq:tame-full-criterion}.
\end{proof}

\subsection{Actual elliptic deformations with an exact tame action}
We next supply the elliptic family and its full-level surface before prescribing a valuation on that surface.

\begin{lemma}\label{tf-lem:surface}
Let $k=\overline{\F}_p$, and let $m$ be one of the allowed tame defects. There is a complete regular local $W(k)$-algebra $R_n$ of dimension two, flat over $W(k)$, carrying an actual good elliptic scheme with a full Drinfeld $p^n$-level structure. It has an integral formal parameter $t$ such that the coordinates $x=t(P)$ and $y=t(Q)$ of a labelled basis form a regular system of parameters. A faithful semilinear $C_m$-action fixes the level labels and $W(k)$ and satisfies
\begin{equation}\label{tf-eq:scalar-action}
 \gamma x=\zeta x,\qquad\gamma y=\zeta y,
\end{equation}
where $\zeta\in W(k)$ has exact order $m$. Its special elliptic action has exact order $m$.
\end{lemma}

\begin{proof}
First choose a supersingular elliptic $E_0/k$ with the required automorphism. For $m=3$ or $6$, the congruence $m\mid p+1$ gives $p\equiv2\pmod3$, and we take $y^2=x^3+1$. For $m=4$, take $y^2=x^3-x$, with $p\equiv3\pmod4$, including $p=3$. The usual scaling automorphisms have the required faithful tangent multipliers. These curves are supersingular: the coefficient of $x^{p-1}$ in the $(p-1)/2$ power of the cubic vanishes, respectively because $3\nmid p-1$ and because $(p-1)/2$ is odd.

For $m=2$, use $[-1]$ on any supersingular elliptic curve. Such a curve exists for every odd $p$ by an elementary Legendre calculation. Put $h=(p-1)/2$. The coefficient just used for $x(x-1)(x-\lambda)$ is $(-1)^h$ times
\[
 H(\lambda)=\sum_{i=0}^h\binom hi^2\lambda^i.
\]
This has positive degree, $H(0)=1$, and
$H(1)=\binom{p-1}{h}\ne0$ in $k$. It has a root distinct from $0,1$, hence a smooth Legendre curve with zero Hasse coefficient. To recall why the coefficient detects supersingularity, write
$y^{-1}=y^{-p}f(x)^{(p-1)/2}$ and apply Cartier to $dx/y$. The degree $3(p-1)/2<2p-1$ shows that only the coefficient of $x^{p-1}$ contributes. Cartier, with its Frobenius twist, is the differential of Verschiebung. Its vanishing makes the degree-$p$ Verschiebung inseparable. Since multiplication by $p$ is Verschiebung composed with the purely inseparable degree-$p$ Frobenius, the $p$-torsion is then connected of rank $p^2$. This is the height-two case.

Let $F_0$ be the formal group of $E_0$. The original deformation theorem of \citet[Propositions~4.2--4.3, pp.~570--573]{Drinfeld1974} gives
$R_0=W(k)[[u]]$ and a finite flat full-level ring $R_n$ which is regular of dimension two, with basis-point coordinates as parameters. The hypotheses allow the integers of any nonarchimedean local field, in particular $\Z_p$ for every odd $p$ \citep[\S1C, pp.~564--565]{Drinfeld1974}. In characteristic zero the level is full: the defining divisor identity gives all $p$-torsion roots, and successive $p$-division gives $p^{2n}$ distinct roots, since multiplication by $p$ is then \emph{\'etale}.

We explain the actual elliptic realization over the complete ring. On each Artinian quotient of $R_0$, the kernels of the universal $[p^j]$ are finite locally free of rank $p^{2j}$ by Weierstrass preparation and form a $p$-divisible group. Apply the Serre--Tate equivalence to this group and $E_0$ \citep[Theorem~1.2.1, p.~143, proof pp.~144--146]{Katz1981}. It gives compatible elliptic deformations. These algebraize directly: the section modules of $\Ocal(j[0])$, $j\ge1$, have their genus-one ranks and compatible reduction maps, since their first cohomology vanishes. Lift bases of the pole-order filtrations to choose compatible Weierstrass coordinates and coefficients. Their inverse limits define a Weierstrass equation over $R_0$ with unit discriminant. The degree-six pole-order relation uses only finitely many coefficients, so this elliptic scheme induces all the Artinian deformations. Its $p$-divisible group is the prescribed one by the compatible finite-level identifications. Pull it back to $R_n$.

The special automorphism $\gamma_0$ acts on the framed deformation functor by changing the residual framing, leaving the level labels fixed. Representability supplies actions on $R_0$ and $R_n$ fixing $W(k)$. Full faithfulness in Serre--Tate supplies the compatible elliptic isomorphisms and their cyclic cocycle identities. These extend to the algebraized Weierstrass model: the finitely many coefficients in the usual Weierstrass coordinate transformation have compatible inverse limits. Thus the action is an actual semilinear action, not only an action on formal isomorphism classes.

Average any parameter $t_0$ by
\[
 t=\frac1m\sum_{i=0}^{m-1}\zeta^{-i}\gamma^i(t_0),
\]
where $\zeta$ is the lift of the special tangent multiplier. Its special derivative is a unit, and $\gamma t=\zeta t$. Changing an integral parameter by a unit linear term preserves the regular system of basis-point parameters. As the labels are fixed, this gives~\eqref{tf-eq:scalar-action}. The special action remains $\gamma_0$ of exact order $m$.
\end{proof}

\subsection{A divisorial valuation and an invariant transverse section}
The next lemma turns the regular surface into a complete field with perfect residue field while retaining its exact index and torsion marking.

\begin{lemma}\label{tf-lem:realization}
Fix $n\ge1$ and an allowed tame defect $m$. For every integer $A\ge1$ with $A\equiv1\pmod m$, there is an elliptic curve with full rational $p^n$-torsion and exact defect $m$ over a complete field with residue $\overline{\F}_p$ and absolute index $E/m$, where
\begin{equation}\label{tf-eq:index-A}
 \frac{E}{p-1}=p^{2n-2-c}(A+p^{2c+1}).
\end{equation}
Here $0\le c<n-1$ when $p^{2c}\le A\le p^{2c+2}$, and $c=n-1$ when $A\ge p^{2n-2}$. Adjacent expressions agree at their common endpoints.
\end{lemma}

\begin{proof}
Use the surface of Lemma~\ref{tf-lem:surface}. Blow up its closed point and then the origins of the successive charts with coordinates $y,x/y^j$, for $1\le j<A$. The final exceptional divisor defines the valuation
\[
 v(y)=1,\qquad v(x)=A,
 \qquad k(v)=k(z),\quad z=x/y^A.
\]
All these charts are regular. For example the first chart is defined in the regular three-dimensional ring $R_n[t]$ by $x-yt$; at a closed point of the exceptional line this relation has a nonzero linear term, leaving regular parameters $y,t-t_0$. The same argument applies successively. The centres are invariant, and on the final chart
\[
 \gamma y=\zeta y,\qquad\gamma z=\zeta^{1-A}z=z.
\]

The minimum and maximum primitive $p^n$-torsion valuations are exactly $1,A$. Here is the cancellation check. A primitive vector $(a,b)$ has coordinate
\[
 f\equiv\bar a x+\bar b y\pmod{(x,y)^2},
 \qquad(\bar a,\bar b)\ne(0,0).
\]
If $A>1$ and $\bar b\ne0$, its value is one. If $A=1$, its leading coefficient is $\bar a z+\bar b\ne0$ in $k(z)$, again giving value one. If $A>1$ and $\bar b=0$, then $\bar a\ne0$. The level-one product factorization of $[p](T)/T$ implies $p\in(x)$: one nonzero $p$-torsion factor is $[p^{n-1}](x)$, itself divisible by $x$. Hence the complete regular quotient $R_n/(x)$ is the characteristic-$p$ DVR $k[[y]]$. Lift the image of $f$ there using Teichm\"uller coefficients to write
\[
 f=ux+g(y),\qquad \bar u=\bar a\ne0,
 \qquad \ord_y g\ge2
\]
with $g=0$ allowed. If $\ord_y g\ne A$, the smaller valuation determines $v(f)\le A$. At equality the leading coefficient is $\bar a z+c\ne0$ in $k(z)$, so $v(f)=A$. Every primitive coordinate has value at least one, and $Q,P$ attain $1,A$. This proves both extrema.

Complete the local DVR at the generic point of the exceptional divisor. It has mixed characteristic because $R_n$ is a domain flat over $W(k)$. The good special formal group remains height two over $k(z)$, and all full-level points remain rational. Lemma~\ref{tf-lem:extrema} applies, despite this intermediate imperfect residue field. Its first equation gives $b=p^{2n-2}$. Its second equation is a continuous strictly increasing piecewise affine function of $M=E/(p-1)$: on
\[
 (p+1)bp^c\le M\le(p+1)bp^{c+1}
\]
it takes the values $p^{2c}\le A\le p^{2c+2}$, with the final interval unbounded. Inverting it gives exactly~\eqref{tf-eq:index-A}. This determines the actual $v(p)$, rather than merely assigning a proposed weight to $p$.

We now specialize transversely. Near the generic point of the exceptional divisor, $p=y^E u$ with $u$ a regular unit after restriction to a nonempty open. Intersect its finitely many translates to make the open invariant. Choose a closed point $z=z_0$ there; only finitely many points of the affine exceptional line are excluded, and $k$ is infinite. Complete the regular local ring at that point and impose
\[
 z-[z_0]=0,
\]
where $[z_0]\in W(k)$ is the Teichm\"uller lift. This is an invariant regular parameter transverse to $y=0$. The quotient $B$ is a complete DVR with uniformizer $y$, residue $k$, and $p=y^E$ times a unit. In particular it has characteristic zero and the same absolute index $E$. The action descends and remains faithful on $y$.

Put $L=\operatorname{Frac}(B)$ and $K=L^{C_m}$. The invariant value group is exactly $m\Z$: the leading residue of an invariant element of valuation $q$ transforms by $\zeta^q$, forcing $m\mid q$, while $y^m$ is invariant. Thus $[L:K]=m$, the fixed ring is a complete DVR of residue $k$, and $e_K=E/m$.

The actual elliptic family pulls back to $B$, and its semilinear generic action descends by finite Galois descent to an elliptic curve over $K$. The labelled full-level sections descend individually. Their generic level remains full because $p\ne0$. The special action has faithful order $m$, so the action on every prime-to-$p$ Tate module has exact order $m$. The good-reduction criterion identifies the exact defect with $m$. Its reduction is additive: it is potentially good and not good, whereas multiplicative reduction would give negative valuation of $j$.
\end{proof}

\subsection{Solving the integer condition and the least-level problem}
We finish by choosing the monomial valuation for each allowed index and then maximizing the torsion level.

\begin{proof}[Proof of Theorem~\ref{thm:tame-full-level}]
Necessity is Proposition~\ref{tf-prop:necessity}. Suppose~\eqref{eq:tame-full-criterion} holds, and put
\[
 M=\frac{me}{p-1},\qquad b=p^{2n-2},\qquad
 c=\min\left\{n-1,\left\lfloor\log_p\frac{M}{(p+1)b}\right\rfloor\right\}.
\]
The bound gives $M\ge(p+1)bp^{\max(0,2n-2-s)}$. Because $s\ge n-1$, the last exponent is at most $n-1$, so
$c\ge\max(0,2n-2-s)$. Hence
\begin{equation}\label{tf-eq:choose-A}
 A=\frac{M}{p^{2n-2-c}}-p^{2c+1}
\end{equation}
is an integer. If $c<n-1$, its defining interval gives
$p^{2c}\le A<p^{2c+2}$; if $c=n-1$, it gives $A\ge p^{2n-2}$. In every case $A\ge1$. Since $m\mid M$, $p\nmid m$ and $p\equiv-1\pmod m$,
\[
 A\equiv-p^{2c+1}\equiv1\pmod m.
\]
Apply Lemma~\ref{tf-lem:realization}. Its good field has absolute index
\[
 E=(p-1)p^{2n-2-c}(A+p^{2c+1})=(p-1)M=me.
\]
Its fixed field therefore has exactly the prescribed index $e$, residue $\overline{\F}_p$, full rational $p^n$-torsion and exact defect $m$. This proves sufficiency for every permitted index.
\end{proof}

\begin{proof}[Proof of Theorem~\ref{thm:tame-level-one}]
At $n=1$, condition~\eqref{eq:tame-full-criterion} reduces to $e\ge(p^2-1)/m$. Lemma~\ref{tf-lem:inertia} and Proposition~\ref{tf-prop:necessity} supply the other necessary conditions. The construction proves their sufficiency. Among $2,3,4,6$, the largest divisor of $p+1$ is respectively $2,6,4,6$ for $p\equiv1,5,7,11\pmod{12}$. This gives the union over all defects. At $p=3$, the two bounds are $e\ge4$ for $m=2$ and $e\ge2$ for $m=4$, with even $e$ in both cases.
\end{proof}

\begin{proof}[Proof of Corollary~\ref{cor:tame-forcing-level}]
Write $m_*=m_{\max}(p)$. If $m_*e<p^2-1$, no nontrivial defect admits full rational $p$-torsion, so the least forcing level is one. Otherwise let $s,L$ be as in the statement. A bad potentially good curve with full rational $p^n$-torsion exists if and only if the criterion holds for $m_*$: all smaller allowable defects give stronger lower bounds, while $m_*$ itself is realized. With $r=n-1$, the exact inequalities are
\[
 r\le s,\qquad 2r\le L,\qquad4r-s\le L.
\]
The largest bad level is therefore
\[
 N=1+\min\left\{s,\left\lfloor L/2\right\rfloor,
                       \left\lfloor(L+s)/4\right\rfloor\right\}.
\]
An actual bad example exists at $N$, and its lower-level torsion is also rational. No bad example exists at $N+1$. For $p\ge5$ all potentially good nontrivial defects were covered by Lemma~\ref{tf-lem:inertia}; for $p=3$ the assertion is restricted to tame inertia and uses $m_*=4$.
\end{proof}
